% =====================================================================
%  capitulos/cap06.tex — Capítulo 6. La estructura de la oración
%  Autora: Mabel Giammatteo
%  Fuente: Capitulo_06_normalizado.docx
% =====================================================================
\capitulo{La estructura de la oración}{Mabel Giammatteo}

\section{Introducción}

Según se ha visto en los capítulos anteriores, la lengua es esencialmente un
mecanismo combinatorio. A partir de elementos mínimos sin significado, como los
fonemas y las sílabas, se van conformando unidades significativas cada vez más
amplias: primero morfemas y, luego, palabras, las que a su vez se proyectan en
unidades intermedias, como los sintagmas (SS). Ahora bien, ¿cómo ubicamos en esta
secuencia de complejidad creciente a la oración? En este capítulo mostramos que
la oración no es solo el resultado de la combinación de un SN, que funciona como
sujeto, y de un SV, que constituye el predicado, sino que también es necesario
dar cuenta de cómo se produce la vinculación entre estos dos constituyentes.
Tanto el tiempo (T), que ancla la oración en un momento dado, como la relación
de concordancia (Conc) entre sujeto y predicado resultan fundamentales para la
constitución de la oración. Es así que ambas categorías ---T y Conc---, reunidas
en la categoría funcional de flexión (F), constituyen el núcleo de la oración,
lo cual convierte a esta unidad en un sintagma de flexión (SF). Se explica
también en qué consiste la relación de predicación y por qué en algunas lenguas
se puede omitir el sujeto y en otras no. Por último, presentamos el concepto de
capa o nivel oracional, que permite dar cuenta de la estructura tripartita de
la oración.

\section{Hacia la definición de oración}\label{sec:06-definicion}

Como se muestra en \exref{ex:06-muchacho}, a pesar de que el contenido puede
resultar más o menos comprensible, la mera yuxtaposición de un SN y un SV no
constituye por sí misma una oración bien formada en español.

\begin{exe}
  \ex[*]{\label{ex:06-muchacho} Muchacho correr.}
\end{exe}

Para explicar la agramaticalidad de \exref{ex:06-muchacho}, podemos, por una
parte, recordar que, para funcionar adecuadamente dentro del SN sujeto de la
oración, el N común necesita estar especificado, es decir que debe ser
antecedido por algún tipo de determinante (D) ---\forma{el / este / tu
muchacho}--- o cuantificador (Q) ---\forma{un / algún / cualquier muchacho}---
(v. \capref{5}, §\ref{sec:05-sn}). Por otra parte, el V necesita manifestar el
tiempo (T) de ocurrencia del evento, como se muestra en \exref{ex:06-corre}.

\begin{exe}
  \ex\label{ex:06-corre} El muchacho corre / corrió / correrá.
\end{exe}

Además, ambos, N y V, deben estar vinculados entre sí mediante la concordancia
(Conc) en número y persona, como muestra la gramaticalidad de
\exref{ex:06-conc-a} frente a oraciones agramaticales como las de
\exref{ex:06-conc-b} y \exref{ex:06-conc-c}.

\begin{exe}
  \ex\label{ex:06-conc}
    \begin{xlist}
      \ex\label{ex:06-conc-a} El muchacho corrió.
      \ex[*]{\label{ex:06-conc-b} El muchacho corrieron.}
      \ex[*]{\label{ex:06-conc-c} Los muchachos corrió.}
    \end{xlist}
\end{exe}

\begin{destacado}
Según vemos, la vinculación entre el SN sujeto y el SV predicado se produce
gracias a las categorías de T y Conc, reunidas en la flexión (F), que se
manifiesta en la desinencia verbal.
\end{destacado}

Aunque se alojan en la desinencia verbal, tanto el T como la Conc trascienden al
verbo. El T es una categoría que afecta a toda la oración; por eso en
\exref{ex:06-conc-a} podemos decir que la oración (y no solo el verbo) está en
pasado.\pendiente{El original remite a (1)} En cuanto a la Conc, es la
categoría que pone en relación al verbo con el sujeto. La F es una categoría
funcional porque, al igual que las palabras funcionales que se estudiaron en el
capítulo anterior, transmite un significado gramatical: la categoría de tiempo
(T) expresa en qué momento ocurre el evento, y la de Conc manifiesta las
categorías de persona y número, que indican quiénes y cuántos intervienen en el
suceso expresado por el verbo.

Pero ¿qué clase de sintagma da como resultado la unión de un SN y un SV?
Recordemos que, según ya se ha visto (v. \capref{5}, §\ref{sec:05-sintagmas}),
todos los sintagmas tienen núcleo y que es el núcleo el que aporta la categoría
a todo el sintagma. Pues bien, dado que la F convierte a la unión de un SN y un
SV en oración (O), al aportarle los rasgos de tiempo (T) y concordancia (Conc),
la categoría funcional F se considera el núcleo de la O.

\begin{destacado}
En consecuencia, la O es un SF (sintagma de flexión) y se define formalmente como
una unidad de construcción que articula un \textbf{SN} y un \textbf{SV} en otro
sintagma de núcleo funcional: el \textbf{SF}.
\end{destacado}

Desde el punto de vista de la estructura sintagmática, el núcleo F toma como su
complemento al SV y aloja en su Esp al SN sujeto (aspecto que precisaremos más
adelante), lo que se puede representar mediante encorchetamiento
\exref{ex:06-sf-a} o con un diagrama arbóreo \exref{ex:06-sf-b}.

\begin{exe}
  \ex\label{ex:06-sf}
    \begin{xlist}
      \ex\label{ex:06-sf-a} {}[\etq{SF/O} [\etq{Esp/SN/Suj} \ldots{}]
          [\etq{F\prima} [\etq{F/n} T, Conc] [\etq{Comp/SV/Pred} \ldots{}]]]
      \ex\label{ex:06-sf-b}
          \begin{forest} gramatica
            [SF / O
              [SN\\(Esp / Suj)]
              [F\prima
                [{F\\T, Conc}]
                [SV\\(Comp / Pred)]]]
          \end{forest}
    \end{xlist}
\end{exe}

\begin{ejercicio}\label{ej:06-determinante}
Compare los sujetos de las oraciones de 1 con los de 2. Proponga alguna
hipótesis para explicar la falta del determinante (D) en los SSNN sujeto de las
oraciones de 2.

\textit{Dato adicional}: tenga en cuenta las subclases del nombre.

\begin{exe}
  \ex
    \begin{xlist}
      \ex Las ratas son las primeras en abandonar el barco.
      \ex Varios niños faltaron hoy a la escuela.
    \end{xlist}
  \ex
    \begin{xlist}
      \ex Pedro llegó cuando nadie lo esperaba.
      \ex Buenos Aires es una ciudad magnífica.
    \end{xlist}
\end{exe}

\textbf{Pista}: en 1, la omisión del D (\juicio{*}{\forma{Ratas son las primeras
en abandonar el barco}}) produce una oración agramatical.
\end{ejercicio}

\begin{ejercicio}\label{ej:06-oracion}
Considere las siguientes oraciones y explique lo pedido más abajo.

\begin{exe}
  \ex Mi padre me prometió un helado.
  \ex Me prometió un helado.
\end{exe}

\begin{consignas}
  \item Fundamente por qué 1 es una oración. Identifique el núcleo de 1 y haga
        la representación entre corchetes.
  \item ¿Considera que 2 también es una oración? ¿Qué particularidad presenta
        respecto de 1?
\end{consignas}

\begin{resuelto}[Punto resuelto]
a) 1 es una O porque contiene un SN sujeto y un SV predicado, vinculados por un
núcleo funcional: la F, que manifiesta T (pasado) y Conc (3.ª pers.~sg.).
Encorchetamiento: [\etq{SF/O} [\etq{SN/Suj} mi\etq{D/Esp} padre\etq{N/n}]
[\etq{F\prima} F\etq{n} [\etq{SV/Pred} me prometió un helado]]].
\end{resuelto}
\end{ejercicio}

\section{La relación sujeto-predicado}\label{sec:06-sujeto}

\subsection{Acerca del sujeto tácito o nulo}\label{sec:06-tacito}

Aunque, según acabamos de plantear, la buena formación de la oración requiere la
combinación de un SN sujeto con un SV predicado, como sabemos, en español el
sujeto puede elidirse, es decir que puede quedar tácito o no estar expresado
léxicamente. En el \capref{1} (v.~§\ref{sec:01-generativa}) vinculamos esta
posibilidad de omitir o no el sujeto que tienen las lenguas con uno de los
principios reconocidos por la gramática generativa: el \termino{principio de
proyección extendida} (PPE), que establece que todas las oraciones deben poseer
un sujeto del cual se predica.\footnote{Este principio es una extensión del
principio de proyección \parencite{chomsky1981}, según el cual la sintaxis debe
respetar los requisitos de selección de las piezas léxicas en todos los niveles
de representación.} Sin embargo, llama la atención que algunas lenguas puedan
elidirlo, según mostramos en \exref{ex:06-esp} y \exref{ex:06-ita}, pero otras
no, como se ejemplifica en \exref{ex:06-ing} y \exref{ex:06-fra}.

\begin{exe}
  \ex\label{ex:06-esp} Español: \forma{Llegaron tarde.}
  \ex\label{ex:06-ita} Italiano: \forma{Erano in ritardo.}
  \ex\label{ex:06-ing} Inglés: \forma{They arrived late.} /
      \juicio{*}{\forma{Arrived late.}}
  \ex\label{ex:06-fra} Francés: \forma{Ils étaient en retard.} /
      \juicio{*}{\forma{Étaient en retard.}}\pendiente{El original dice
      «étaint»}
\end{exe}

¿Cómo podemos explicar esta diferencia? Según ya hemos visto en el \capref{1},
los principios de carácter universal se realizan mediante parámetros que, por lo
general, presentan opciones binarias, lo cual significa que existen dos
posibilidades de realización, en este caso, del sujeto: la presencia obligatoria
en todas las oraciones, como sucede en inglés y en francés, frente a la
posibilidad facultativa de omitirlo, que se da en lenguas como el español y el
italiano, entre otras. Esta segunda opción se conoce como \termino{parámetro del
sujeto nulo} y las lenguas que permiten la omisión del sujeto se conocen como
\textit{pro-drop} (con caída del pronombre).

Ahora bien, ¿por qué en algunas lenguas es posible la omisión, mientras que
otras la rechazan? Las lenguas \textit{pro-drop} tienen una morfología verbal
rica, que permite reponer el sujeto a partir de las categorías morfológicas de
la desinencia verbal.\footnote{No obstante, hay lenguas como el alemán, con una
rica morfología, que no permiten la omisión, y otras como el chino, que no
expresan categorías de concordancia e igualmente permiten la omisión. No nos
extenderemos aquí sobre estas particularidades, que exceden el nivel inicial de
este texto.} Así, \forma{llegaron} \exref{ex:06-esp} y \forma{erano}
\exref{ex:06-ita} remiten a una tercera persona plural (\forma{ellos / ellas},
para el español, y \forma{loro}, femenino y masculino, para el italiano).

Por lo anteriormente expuesto, en la representación estructural de estas
oraciones se incluye un \textit{pro} (= pronombre) en la posición de sujeto,
como se muestra en \exref{ex:06-pro-a} y se representa en el diagrama arbóreo
\exref{ex:06-pro-b}.\footnote{El uso del triángulo en el esquema arbóreo
indica que, a fin de simplificar la descripción, se omite el análisis interno de
un constituyente. En el caso de \exref{ex:06-pro-b}, no se realiza el análisis
interno del SV \forma{llegaron tarde}.}

\begin{exe}
  \ex\label{ex:06-pro}
    \begin{xlist}
      \ex\label{ex:06-pro-a} {}[\etq{SF/O} \textit{pro}\etq{Esp/Suj}
          [\etq{F\prima} F\etq{n} [\etq{SV/Pred/Comp} llegaron tarde]]]
      \ex\label{ex:06-pro-b}
          \begin{forest} gramatica
            [SF / O
              [\textit{pro}]
              [F\prima
                [F]
                [SV [\textit{llegaron tarde}, roof]]]]
          \end{forest}
    \end{xlist}
\end{exe}

\begin{recuadro}[Para ampliar: ¿qué clase de pronombre es \textit{pro}?]
El \textit{pro} que hace las veces de sujeto en las oraciones de sujeto tácito se
caracteriza como una \termino{categoría vacía} o \termino{nula} que permite que
el PPE se respete. Se trata de un elemento no audible (sin rasgos fonológicos),
pero estructuralmente necesario para rellenar la posición de sujeto. Tal como se
muestra en \exref{ex:06-pro-b}, al igual que los sujetos expresos, \textit{pro}
se legitima en la posición de Esp de la F y se considera que sus rasgos
gramaticales de persona y número concuerdan con los de la desinencia verbal.
Una característica de \textit{pro} es que, debido a que carece de rasgos
fonológicos, no puede focalizarse (v.~§\ref{sec:06-capas}) o ponerse de relieve,
ya que este procedimiento requiere que el elemento destacado reciba acento
fuerte y, obviamente, un elemento que no se pronuncia no puede ser acentuado.
Así, mientras el sujeto léxico de \exref{ex:06-foco-a} puede destacarse mediante
la focalización, el sujeto tácito de \exref{ex:06-foco-b} no puede realzarse
entonativamente.

\begin{exe}
  \ex
    \begin{xlist}
      \ex\label{ex:06-foco-a} Ganó la beca ELLA, no su hermano.
      \ex[*]{\label{ex:06-foco-b} Ganó la beca ∅, no su hermano.}
    \end{xlist}
\end{exe}
\end{recuadro}

\subsubsection{La interpretación del sujeto tácito}

Según \textcite{bosque2009}, el contenido del sujeto tácito \textit{pro} se
recupera a partir del discurso precedente \exref{ex:06-andes1}, del
subsiguiente \exref{ex:06-andes2} o a partir del contexto extraoracional
\exref{ex:06-andes3}:

\begin{exe}
  \ex\label{ex:06-andes1} Los pueblos originarios de los Andes\etq{i} no
      conocían el caballo, aunque \textit{pro}\etq{i}\footnote{El subíndice que
      sigue a \textit{pro} indica cuál es el elemento del contexto al que
      \textit{pro} remite, el cual también va seguido de dicho índice; por eso
      se dice que ambos elementos están coindizados.} utilizaban llamas como
      animales de carga.
  \ex\label{ex:06-andes2} Aunque \textit{pro}\etq{i} no conocían el caballo, los
      pueblos originarios de los Andes\etq{i} utilizaban llamas como animales de
      carga.
  \ex\label{ex:06-andes3} \textit{pro}\etq{[1.ª pers. sg.]} Leí que los pueblos
      originarios utilizaban llamas como animales de carga.
\end{exe}

\subsection{Sujetos tácitos vs. sujetos explícitos}\label{sec:06-explicitos}

Para ver cuáles son las condiciones en que se usan los sujetos tácitos,
consideremos comparativamente las siguientes secuencias de oraciones:

\begin{exe}
  \ex\label{ex:06-nino} Descubrieron a un niño perdido en la montaña. Tenía las
      manos congeladas. Estaba muy asustado y no podía expresarse bien.
  \ex\label{ex:06-llegaron} Llegaron de noche. Estaban sumamente fatigados. Nadie
      los vio entrar. Traían las manos heladas y estaban muy desabrigados.
      Habían venido huyendo por leguas y, sin embargo, aún no tenían un destino
      seguro.
\end{exe}

Mientras que en \exref{ex:06-nino}, a partir de su mención en la primera
oración ---\forma{un niño perdido en la montaña}---, podemos fácilmente
recuperar el sujeto de la segunda y tercera oración, en \exref{ex:06-llegaron},
en cambio, aunque entendemos perfectamente lo que dice el texto, y también
podemos identificar que se está haciendo referencia a una tercera persona
plural, no podemos inferir de quiénes se está hablando en el fragmento. Tal como
plantea \textcite{camacho2018}, la interpretación de los sujetos nulos requiere
un \termino{tópico}\footnote{Aquí utilizamos \textit{tópico} o \textit{tema
discursivo} en el sentido general referido a aquello de lo que se está
hablando.} ya introducido en el discurso a partir del cual estos puedan ser
identificados. Además, los sujetos nulos también se emplean cuando el tópico
discursivo se mantiene; así, en \exref{ex:06-nino}, el tópico \forma{un niño
perdido en la montaña} permanece como sujeto en las dos oraciones siguientes.

Y ¿cuándo se utilizan los sujetos explícitos? Veamos lo que sucede en
\exref{ex:06-osos}.

\begin{exe}
  \ex\label{ex:06-osos} Una fotógrafa observó el comportamiento de los osos
      polares. Una cría de oso polar debe aprender las habilidades de
      supervivencia necesarias en el Ártico. Las madres se encargan de proteger
      a sus crías. Los machos, en cambio, no se hacen cargo del cuidado de los
      oseznos.
\end{exe}

El ejemplo de \exref{ex:06-osos} muestra que los sujetos explícitos se usan en
el comienzo de un texto o párrafo para identificar al referente del que estamos
hablando ---\forma{una fotógrafa}, en la primera oración--- y también se emplean
cuando se produce un cambio de tópico; así, del sujeto \forma{una fotógrafa} en
la primera oración se pasa a \forma{una cría de oso polar}, en la segunda; a
\forma{las madres}, en la tercera, y a \forma{los machos}, en la última.

\begin{ejercicio}\label{ej:06-peyrou}
Lea atentamente el siguiente fragmento adaptado de \textit{Se vuelven contra
nosotros}, de Manuel Peyrou, y luego realice las actividades pedidas, según se
muestra en los ejemplos resueltos.

\begin{consignas}
  \item Encorchete los sujetos de todas las oraciones y subraye los verbos a los
        que se vinculan.
  \item Inserte un \textit{pro} en el lugar correspondiente a los sujetos nulos y
        señale sus rasgos de concordancia (número y persona).
  \item Cuando sea posible, identifique el antecedente de cada sujeto nulo y
        proponga alguna hipótesis explicativa para aquellos casos en que no sea
        posible reconocer un antecedente en el texto.
\end{consignas}

% Fragmento literario protegido: se incorpora a mano desde el original.
\textoaincorporar{el fragmento adaptado de \textit{Se vuelven contra nosotros},
de M.~Peyrou, con las oraciones numeradas de (1) a (8): desde «Estos hombres
llevaban el nombre genérico de \textit{Proveedores}» hasta «\ldots{}y el hombre
dobló su fortuna en tres meses»}\pendiente{El original fecha el texto en 2020 y
escribe «Peyroú»}

\begin{resuelto}
a) (1) [\etq{Suj} Estos hombres]\etq{i} llevaban el nombre genérico de
«Proveedores». b) (2) A veces realmente [\etq{Suj}
\textit{pro}\etq{[3.ª pl.]}]\etq{i} vendían caballos y otros elementos
necesarios. c) El antecedente de \textit{pro} es \forma{estos hombres}, de la
oración anterior.
\end{resuelto}
\end{ejercicio}

\begin{recuadro}[Para ampliar: un caso de variación dialectal en el uso de los sujetos explícitos]
Aunque en el español general la reiteración del sujeto cuando no hay cambio de
tópico resulta redundante, en el español de la región de El Cibao, en República
Dominicana, los sujetos pronominales tienden a repetirse.

\begin{exe}
  \ex Cuando \textbf{tú} acabe \textbf{tú} me avisas.
      \parencite[ejemplo tomado de][393]{rodriguezramalle2005}
\end{exe}

Para Rodríguez Ramalle, esto sucede porque las marcas flexivas, que en el
español general permiten omitir el sujeto, se debilitan hasta no pronunciarse a
veces, lo que provoca la necesidad de reposición del sujeto.
\textcite{bullock2009} plantean que en esta variedad los pronombres explícitos de
sujeto son mucho más frecuentes que en el español general e incluso se usan,
contrariamente a lo esperado, para continuar el tópico discursivo, como muestra
el siguiente ejemplo tomado de los citados autores:

\begin{exe}
  \ex Entonces \textbf{él} fue a buscar el gallo, en busca el resto de las
      plumas, y \textbf{él} fue a buscarlo. Entonces, \textbf{él} vio un oso
      debajo de la mata de mango y vio el gallo a la parte de atrás del oso.
      \parencite[56; hablante dominicano de El Cibao]{bullock2009}
\end{exe}
\end{recuadro}

\begin{recuadro}[Diferencia entre oración y enunciado]
Según hemos planteado, la oración se define fundamentalmente como una unidad
propia de la gramática, formada por un SN sujeto y un SV predicado, articulados
por la F; se trata, pues, de una unidad estructural teórica, independientemente
de su empleo en situaciones concretas de habla. Para las emisiones efectivamente
proferidas en una situación concreta por un hablante determinado y dirigidas a
un oyente también determinado, se emplea la denominación de
\termino{enunciados}. Por tanto, mientras la oración es una unidad de la
gramática, el enunciado es una unidad del uso que no requiere tener la
estructura predicativa. Así, unidades como \forma{¡Socorro!, En la otra esquina,
Bien} o \forma{De ninguna manera} son enunciados perfectamente comprensibles y
que podríamos emplear en situaciones comunicativas concretas, aunque no pueda
reconocerse en ellos la estructura predicativa de sujeto y predicado.
\end{recuadro}

\section{Predicados no verbales y predicaciones sin sujeto}\label{sec:06-noverbales}

Para la \textit{Nueva gramática de la lengua española} \parencite[§\,1.13a]{rae2025},
las oraciones «son unidades mínimas de predicación organizadas en torno a un
verbo, que ponen en relación un sujeto con un predicado». El texto aclara que el
predicado puede ser verbal o no \parencite[§\,1.12d]{rae2025}. Al respecto
podemos tomar un caso como \forma{Interesante, la película}, en el que, aunque no
hay verbo, igualmente hay predicación. Se trata de una predicación de tipo
enfático, que se marca no solo por la ausencia del verbo cópula, que se podría
reponer ---\forma{Interesante es la película}---, sino por la alteración del
orden entre sujeto y predicado, que se vuelve significativa. En este caso
tenemos un predicado no verbal ---\forma{interesante}---, referido a un sujeto
pospuesto ---\forma{la película}---, junto al cual constituye una estructura
oracional de tipo enfático.

Ahora bien, en casos como \forma{Llueve} o \forma{Hace frío}, tenemos un verbo
conjugado, pero, dado que se trata de eventos climáticos, de naturaleza
impersonal, no hay un sujeto referencial con el que se los pueda relacionar. No
obstante, se trata de oraciones porque las formas verbales conjugadas
\forma{llueve} y \forma{hace (frío)} gramaticalmente manifiestan tercera persona
singular. Nuevamente, para el español, tendremos un sujeto \textit{pro}, ya que
en las lenguas \textit{pro-drop} los verbos de este tipo no pueden tener un
sujeto expreso. Y ¿qué sucede en estos casos con las lenguas que realizan la
primera opción del parámetro del sujeto nulo, como el inglés y el francés? Como
en estas lenguas la presencia de un sujeto es obligatoria, la posición de sujeto
es ocupada por un \termino{expletivo}, es decir, un elemento con rasgos
fonológicos, pero vacío de significado, como \forma{it} en inglés
\exref{ex:06-it} o \forma{il} en francés \exref{ex:06-il}.

\begin{exe}
  \ex\label{ex:06-it} \forma{It rains.} (\sig{Llueve})
  \ex\label{ex:06-il} \forma{Il pleut.} (\sig{Llueve})
\end{exe}

Las lenguas \textit{pro-drop}, en cambio, no necesitan un expletivo léxico, dado
que el pronombre sujeto puede ser tácito y estar representado por un
\textit{pro}\etq{expl}, tal como sucede en español \exref{ex:06-llueve} y en
italiano \exref{ex:06-piove}.

\begin{exe}
  \ex\label{ex:06-llueve} {}[\etq{SF/O} [\etq{Suj} \textit{pro}\etq{expl}]
      [\etq{SV/Pred} llueve\etq{V/n}]]
  \ex\label{ex:06-piove} {}[\etq{SF/O} [\etq{Suj} \textit{pro}\etq{expl}]
      [\etq{SV/Pred} piove\etq{V/n}]]
\end{exe}

En español también existe otro tipo de oraciones impersonales, con verbo en
tercera persona plural, como las siguientes:

\begin{exe}
  \ex\label{ex:06-arb}
    \begin{xlist}
      \ex\label{ex:06-arb-a} Te esperan en la puerta.
      \ex\label{ex:06-arb-b} Me llamaron por teléfono hoy temprano.
    \end{xlist}
\end{exe}

En oraciones como \exref{ex:06-arb-a} y \exref{ex:06-arb-b} podemos interpretar
que la referencia de la flexión remite a un sujeto plural mencionado
anteriormente y recuperable del contexto, tal como se planteó en
§\ref{sec:06-explicitos} ---por ejemplo, \forma{Vinieron tus amigos a buscarte.
Te esperan en la puerta}---. En este caso el sujeto plural corresponde a
entidades definidas e identificables: \forma{ellos / tus amigos}. Pero los
ejemplos de \exref{ex:06-arb} también admiten otra interpretación, muy común en
la lengua, en la que no hay una referencia a una entidad determinada y en la que
«el plural es una mera marca de indefinición, ya que puede referirse a un único
individuo» \parencite[109]{ditullio2014}. Este sujeto, que se interpreta como
indefinido e inespecífico, se conoce como \termino{pro arbitrario}:

\begin{exe}
  \ex
    \begin{xlist}
      \ex {}[\etq{SF/O} [\etq{Suj} \textit{pro}\etq{arb}] [\etq{SV/Pred} te
          esperan en la puerta]]
      \ex {}[\etq{SF/O} [\etq{Suj} \textit{pro}\etq{arb}] [\etq{SV/Pred} me
          llamaron por teléfono hoy temprano]]\pendiente{El original dice «Te
          llamaron»}
    \end{xlist}
\end{exe}

\begin{ejercicio}\label{ej:06-esoracion}
Decida cuáles de las siguientes son oraciones y fundamente su respuesta, según
se muestra en los ejemplos resueltos de abajo.

\begin{exe}
  \ex Mario es el barbero del barrio.
  \ex Un sinfín de problemas.
  \ex Cayó un meteorito en la autopista.
  \ex Vendrán más tarde.
  \ex Allá en el rancho grande.
  \ex Allá, solo ranchos y casuchas con techo de paja.
  \ex En el cuento había un ogro, un burro y una princesa.
  \ex Primero entraron los familiares cercanos y después, los demás parientes.
  \ex Granizó copiosamente en varias zonas del interior de la provincia.
  \ex Infaltable, el café de la sobremesa.
\end{exe}

\begin{resuelto}[Ejemplos resueltos]
1. O formada por el SN sujeto \forma{Mario} y el SV predicado \forma{es el
barbero del barrio}. 2. No es O, sino SN.
\end{resuelto}
\end{ejercicio}

\subsection{La relación predicativa}\label{sec:06-predicativa}

Ahora que ya hemos presentado la estructura oracional de acuerdo con el esquema
básico y universal de la teoría X\prima{} (v. \capref{5},
§\ref{sec:05-xbarra}), podemos también preguntarnos qué significa que el SN y el
SV se combinen para formar una oración. La relación que se establece entre los
dos constituyentes esenciales de la oración, el SN y el SV, se conoce como
\termino{predicación}. En ella el verbo asume un doble papel: por un lado, como
categoría léxica, se desempeña como núcleo del SV, soporte del predicado; por el
otro, en su desinencia, el verbo aloja a la flexión, categoría funcional que,
mediante la concordancia y la temporalización, instaura la relación
predicativa.

En tanto relación predicativa, toda oración manifiesta un evento o estado y, por
lo tanto, se organiza en torno al verbo que lo expresa. La mayor parte de los
verbos no se emplean solos, sino que tienen distintas exigencias de
complementación sintáctico-semánticas, que son satisfechas por los argumentos
que los acompañan. Como veremos más adelante, los argumentos del verbo remiten a
entidades del mundo extralingüístico que participan en el evento descripto por
el verbo. Por lo tanto, según su semántica, los verbos pueden requerir más o
menos argumentos. Así, por ejemplo, \forma{trabajar} solo exige un argumento,
referido a \sig{la persona que trabaja}\etq{x} \exref{ex:06-trabajar}; en
cambio, \forma{escribir} requiere dos argumentos: \sig{el que escribe}\etq{x} y
\sig{lo que se escribe}\etq{y} \exref{ex:06-escribir}; y hay también verbos que
piden tres argumentos, como, por ejemplo, \forma{enviar}, que debe manifestar
\sig{el que envía}\etq{x}, \sig{lo que se envía}\etq{y} y \sig{a quien se
envía}\etq{z} \exref{ex:06-enviar} (v. \capref{7}).

\begin{exe}
  \ex\label{ex:06-trabajar} \textit{trabajar}(x). Ej.: [Pedro]\etq{x} trabaja.
  \ex\label{ex:06-escribir} \textit{escribir}(x, y). Ej.: [Pedro]\etq{x}
      escribe [un correo electrónico]\etq{y}.
  \ex\label{ex:06-enviar} \textit{enviar}(x, y, z). Ej.: [Pedro]\etq{x} envía
      [un correo electrónico]\etq{y} [a sus alumnos]\etq{z}.
\end{exe}

\begin{ejercicio}\label{ej:06-argumentos}
Proporcione los ejemplos solicitados, según se muestra en el ejemplo de abajo.

\begin{exe}
  \ex Dos oraciones con verbos que exijan un solo argumento.
  \ex Dos oraciones con verbos que exijan dos argumentos.
  \ex Dos oraciones con verbos que exijan tres argumentos.
\end{exe}

\begin{resuelto}[Respuesta posible]
1. a) [Juan] trabaja en el bar de la otra cuadra. El V \forma{trabajar} exige un
solo argumento: \forma{Juan}. \forma{En el bar de la otra cuadra} es un
circunstancial no exigido por el V. b) [\etq{Suj} La pobre muchacha] bailó
frenéticamente hasta casi desfallecer. El V \forma{bailar} exige un solo
argumento. Tanto \forma{frenéticamente} como \forma{hasta casi desfallecer} son
circunstanciales no exigidos por el V.
\end{resuelto}
\end{ejercicio}

\begin{ejercicio}\label{ej:06-funciones}
Considere las siguientes oraciones y luego responda lo pedido, como se muestra
para el ejemplo resuelto~1.

\begin{exe}
  \ex La abogada me entregó los papeles en su oficina.
  \ex El primer ministro informó a la población sobre las nuevas medidas de
      gobierno.
  \ex El fiscal del caso guardó los expedientes en un sobre.
\end{exe}

\begin{consignas}
  \item ¿Cuántos argumentos exigen los verbos de las oraciones anteriores?
  \item ¿Qué funciones sintácticas cumplen los argumentos
        identificados?\footnote{Para las funciones sintácticas, puede consultar
        el \capref{7}.}
\end{consignas}

\begin{resuelto}
1. a) \forma{Entregó} exige tres argumentos: \forma{la abogada, los
certificados} y \forma{el nuevo dueño}. b) [\etq{Suj} La abogada] entregó
[\etq{COD} los certificados] [\etq{COI} al nuevo dueño].\pendiente{La resolución
no corresponde a la oración 1: «La abogada me entregó los papeles en su
oficina»}
\end{resuelto}
\end{ejercicio}

\begin{ejercicio}\label{ej:06-llover}
Observe la siguiente oración del español y responda lo pedido más abajo:

\begin{center}
\forma{Llovió torrencialmente en el fin de semana.}
\end{center}

\begin{consignas}
  \item ¿Cuántos argumentos exige \forma{llover} en la oración anterior?
  \item ¿Se le ocurren otros verbos que se construyan como \forma{llover}?
        Proporcione dos ejemplos oracionales.
  \item ¿Qué sucede en lenguas como el inglés o el francés con los verbos
        identificados en el punto anterior?
  \item ¿Qué nombre recibe el elemento que hace las veces de sujeto en estas
        lenguas?
  \item ¿Con qué parámetro lingüístico puede relacionar la diferencia existente
        en estas lenguas respecto de otras como el español y el italiano?
  \item ¿Qué nombre reciben estas dos últimas lenguas en función de esta
        característica?
\end{consignas}

\begin{resuelto}[Punto resuelto]
a) \forma{Llover} es un verbo impersonal que no exige ningún argumento, ni
siquiera sujeto. \forma{Torrencialmente} y \forma{en el fin de semana} son
circunstanciales no exigidos.
\end{resuelto}
\end{ejercicio}

\begin{ejercicio}\label{ej:06-cuatro}
Algunos autores, como \textcite[268]{bosque2009}, plantean la posibilidad de
verbos con hasta cuatro argumentos y ofrecen ejemplos como los siguientes:

\begin{exe}
  \ex Juan viaja desde Acassuso a Mar del Plata por la ruta dos.
  \ex María le compró su auto a su prima en 500\,000\,\$.
\end{exe}

\begin{consignas}
  \item Encorchete los argumentos de las oraciones anteriores, como se muestra
        en el ejemplo resuelto~1.
  \item Proporcione dos ejemplos más de verbos que admitan cuatro argumentos.
\end{consignas}

\begin{resuelto}
1. a) [Juan] viaja [desde Acassuso] [a Mar del Plata] [por la ruta dos].
\end{resuelto}
\end{ejercicio}

\subsection{Diferencias entre argumentos: argumento externo y argumentos internos}\label{sec:06-externo}

No todos los argumentos que exige el verbo están en el mismo nivel. Por
ejemplo, si bien un verbo transitivo como \forma{construir} exige dos argumentos
que dan cuenta, respectivamente, de \sig{el que construye} y de \sig{lo
construido} \exref{ex:06-arquitecto}, existe una diferencia entre ellos que lleva
a la distinción entre argumento externo y argumento interno
\parencite{williams1980}. El \termino{argumento externo}, que se identifica con
el sujeto por su relación estructural prominente, es el que mantiene con el
verbo la relación de predicación; en cambio, los \termino{argumentos internos} o
complementos verbales tienen una vinculación más estrecha con el verbo y, según
veremos más adelante, permanecen dentro del SV.

\begin{exe}
  \ex\label{ex:06-arquitecto} {}[\etq{Arg. ext.} El arquitecto] construyó
      [\etq{Arg. int.} un puente].
\end{exe}

\begin{recuadro}[Para ampliar: una relación muy estrecha]
Los argumentos internos, especialmente el COD, completan o complementan el
significado del V. Así, aunque se trate del mismo verbo, el significado de
\forma{servir} no es exactamente el mismo en cada una de las siguientes
oraciones, y esa variación se debe a la diferente naturaleza de los objetos que
lo acompañan en cada una de ellas:

\begin{exe}
  \ex El mozo \textit{sirvió} a los invitados (= \sig{los atendió}).
  \ex Ese hombre \textit{sirve} a Dios desde su juventud (= \sig{está consagrado
      a su servicio}).
  \ex Trajeron un toro campeón para \textit{servir} a las vacas (=
      \sig{fecundarlas, prestar servicios sexuales}).
\end{exe}

Por tanto, según estamos viendo, la vinculación entre el V y su objeto es
siempre más estrecha que la que existe entre el V y su sujeto. En español, hay
dos tipos de formaciones que muestran la integración entre el V y su OD, que los
lleva a constituir una unidad: los compuestos V + N y las locuciones verbales.
Dichas formaciones solo se dan entre el V y su OD y nunca entre el V y su
sujeto.

\begin{itemize}
  \item Los compuestos V + N del tipo de \forma{sacacorchos, abrelatas,
        sacapuntas, tirabombas, matamoscas, vendehumo, rompecorazones} (v.
        \capref{4}, §\ref{sec:04-composicion}).
  \item Diversas locuciones verbales, como \forma{hacer la cama} (=
        \sig{engañar}), \forma{dar fe} (= \sig{atestiguar}), \forma{echar
        raíces} (= \sig{radicarse}), \forma{escupir el asado} (=
        \sig{estropear}) (v. \capref{11}).
\end{itemize}
\end{recuadro}

\section{La organización tripartita de la oración: las capas oracionales}\label{sec:06-capas}

En tanto relación predicativa, toda oración manifiesta un evento o suceso y,
por lo tanto, se organiza en torno al verbo que lo expresa. En el denominado
\termino{enfoque cartográfico} \parencite{rizzi1995} se propone que la oración
se estructura en tres capas o niveles (\textit{shells}):

\begin{itemize}
  \item la \termino{capa léxica} o del SV, donde se produce la vinculación del V
        con sus argumentos (incluido el sujeto);
  \item la \termino{capa flexiva}, oracional o del SF, que pone en relación el
        sujeto con el predicado y en la que la oración se sitúa temporalmente;
  \item la \termino{capa informativa}, complementizadora o del SComp, que
        incluye la perspectiva del hablante y establece la articulación con el
        nivel del discurso.
\end{itemize}

Tal como sostiene \textcite{carnie2008}, se trata de

\begin{cita}
uno de los descubrimientos más interesantes de la lingüística del siglo veinte:
[ha sido] la naturaleza en tres capas de las oraciones. Con esto me refiero a la
forma en que un cierto número de enfoques sobre la estructura sintáctica han
convergido en la idea de que, translingüísticamente, las oraciones tienen en su
centro una predicación nuclear, que está dentro de una estructura de operadores
que ubican la predicación en tiempo y espacio respecto del tiempo de la emisión,
y finalmente están encerradas por una estructura que relaciona la proposición
con la perspectiva del hablante. (Traducción nuestra)
\end{cita}

Es importante destacar que esta organización abstracta se refleja en la forma
en que está estructurada la oración: desde el interior hacia el exterior, o sea,
de derecha a izquierda.

\begin{exe}
  \ex {}[\etq{SC} Quizás [\etq{SF} el profesor [\etq{SV} nos dé tarea para el fin
      de semana]]]
\end{exe}

Desde esta perspectiva, la capa léxica o del SV, que se organiza en torno al
verbo (aún no temporalizado), es la más interna. En ella la predicación se
despliega como un evento y el verbo, elemento central de la estructura
oracional, se relaciona con sus argumentos y con los modificadores no exigidos o
adjuntos.

\begin{exe}
  \ex\label{ex:06-castores} {}[\etq{SV} construir [\etq{Arg.~1} diques]
      [\etq{Arg.~2} los castores] [\etq{Adj} en ríos y arroyos]]
\end{exe}

Como se puede advertir en \exref{ex:06-castores}, \forma{construir} manifiesta
un evento de \sig{acción} (v. \capref{8}), que exige dos argumentos: \sig{lo
construido} ---\forma{diques}--- y \sig{quien construye} ---\forma{los
castores}---. La descripción del evento incluye un modificador no exigido por el
V, el adjunto de lugar ---\forma{en ríos y arroyos}---.

El siguiente nivel estructural se conoce como \termino{capa flexiva} o de la
oración porque, como ya adelantamos, mediante la flexión verbal ---que expresa T
y Conc--- el evento queda localizado respecto del momento de la emisión y se
establece la vinculación entre el SV y su sujeto. Es en este nivel en el que el
argumento sujeto debe desplazarse desde la posición interna que como argumento
del V ocupa dentro del SV, como se muestra en \exref{ex:06-viaje}, hacia el Esp
de la F, lo cual lo convierte en un «argumento externo», en tanto los demás
argumentos, que permanecen en el ámbito del SV, se identifican, según ya hemos
dicho, como «argumentos internos».

\begin{recuadro}[Para ampliar: el «viaje» del sujeto]
Llegados a este punto, ya podemos preguntarnos por qué se produce este
movimiento del sujeto hacia el exterior del SV, es decir, qué lo motiva. Además
de la temporalización y la concordancia, la conformación de la oración también
requiere que todos los argumentos del verbo «se vuelvan visibles» o manifiestos
para la sintaxis, para lo cual deben tener una marca que indique la función
sintáctica que cumplen. Ahora bien, ¿cuál es la categoría gramatical que marca
la función sintáctica de los SSNN? Según se vio en el \capref{4}
(v.~§\ref{sec:04-categorias}), el \termino{caso} es la categoría gramatical que
tiene esta función. La marca de caso «vuelve visibles» para su funcionamiento
oracional a los SSNN presentes en la oración. Recordemos que en las lenguas
clásicas ---griego y latín--- y también en algunas lenguas modernas, como el ruso
o el alemán, el caso es una categoría morfológica incluida en la desinencia de
la palabra: \forma{capra} (\sig{cabra}, caso nominativo o del sujeto);
\forma{capram} (\sig{a la cabra}, caso acusativo o del COD). En lenguas como el
español, en las que no existe el caso morfológico, se dice que esta marca es
abstracta, excepto en los pronombres personales, que varían en caso. Así, la
forma del pronombre de primera persona singular para la función de sujeto es
\forma{yo} ---\forma{\textbf{Yo} canto muy bien}---; para la función de objeto
es \forma{me} ---\forma{\textbf{Me} regalaron un hermoso cachorrito}---, y la
forma que se usa como complemento de preposición es \forma{mí} ---\forma{Lo
quiero para \textbf{mí}}---. En la sintaxis, el caso que le corresponde a un SN
determinado debe ser asignado. Así, por ejemplo, el V es la categoría léxica que
asigna caso acusativo al COD. El OI, en cambio, recibe su marca de caso de la
preposición \forma{a} que siempre lo precede: \forma{Les traje regalos
\textbf{a} los niños}. Y ¿qué sucede, entonces, con el sujeto? ¿Cómo recibe
caso? Para recibir caso nominativo, que es el que le corresponde, el sujeto,
originariamente un argumento interno a la capa léxica, se desplaza del SV al Esp
de la F. Por tanto, el sujeto recibe caso nominativo de una categoría funcional:
la F. Asimismo, el V también se desplaza del SV al núcleo de la F para cotejar
sus rasgos flexivos de tiempo y concordancia con el sujeto.

\begin{exe}
  \ex\label{ex:06-viaje}
    \begin{xlist}
      \ex {}[\etq{SF} [\etq{SN1} los castores]\etq{i} [\etq{F\prima} F
          [\etq{SV} [\etq{SN1} \sout{los castores}]\etq{i} [\etq{V\prima}
          construyen\etq{V} [\etq{SN2} diques]]]]]
      \ex
          \begin{forest} gramatica
            [SF
              [SN\etq{1}\\\textit{los castores}\etq{i}, name=suj]
              [F\prima
                [F]
                [SV
                  [SN\etq{1}\\{\color{gristexto}\textit{los castores}}\etq{i}, name=huella]
                  [V\prima
                    [V\\\textit{construyen}]
                    [SN\etq{2}\\\textit{diques}]]]]]
            \draw[->, densely dashed] (huella.south west) to[out=200, in=250]
              (suj.south);
          \end{forest}
    \end{xlist}
\end{exe}
\end{recuadro}

Dado que la oración ya queda conformada mediante la vinculación del SN sujeto y
del SV predicado dentro de la capa flexiva, resulta lícito preguntarnos para qué
se necesita proponer un nivel más o, mejor dicho, ¿qué aporta el tercer nivel
oracional? Este tercer nivel, al que se lo conoce como \termino{capa
informativa} o del \termino{complementante} (C), es el lugar en el que la
oración se conecta con la perspectiva del hablante, por lo que se ubican allí
distintos elementos que expresan la fuerza ilocutiva de la oración, tales como
los adverbios modales, el foco y el tópico, entre otros. Vayamos viendo de qué
se trata esto.

El C es, al igual que la F, un núcleo funcional y, así como la F contiene
valores de T y Conc, el C es de naturaleza modal. Si bien no vamos a
extendernos aquí sobre el complementante (v. capítulos 12 y 13), podemos
ilustrar su carácter modal mostrando la diferencia entre \exref{ex:06-que} y
\exref{ex:06-si}.

\begin{exe}
  \ex
    \begin{xlist}
      \ex\label{ex:06-que} Dice [\etq{SC} \textbf{que} [\etq{SF/O} Juan [\etq{SV}
          viene temprano]]].
      \ex\label{ex:06-si} Pregunta [\etq{SC} \textbf{si} [\etq{SF/O} Juan
          [\etq{SV} viene temprano]]].
    \end{xlist}
\end{exe}

En ambas oraciones, el complemento del verbo matriz o principal ---\forma{dice,
pregunta}--- no es un SN, como en \forma{Dice [la verdad]}, sino una oración
(SF) ---[\forma{Juan viene temprano}]---, por lo que la vinculación entre la
oración principal y la subordinada no es directa, sino mediada por una
conjunción, que puede ser \forma{que} o \forma{si}.

Según vemos, la distribución de ambas conjunciones ---\forma{que} y \forma{si}---
no es aleatoria. Pero ¿a qué se debe la alternancia? Sin duda no depende de la
oración subordinada, ya que en ambos casos ---\exref{ex:06-que} y
\exref{ex:06-si}--- es la misma ---\forma{Juan viene temprano}---. La diferencia
radica, entonces, en el verbo matriz: \forma{decir} en \exref{ex:06-que} y
\forma{preguntar} en \exref{ex:06-si}. Por tanto, es la modalidad de la oración
matriz, declarativa en \exref{ex:06-que} e interrogativa en \exref{ex:06-si}, la
que provoca la diferencia.

\begin{destacado}
\textbf{Una distribución complementaria}
\begin{itemize}
  \item Cuando el verbo de la oración principal es declarativo ---por ejemplo,
        \forma{decir, informar, afirmar}---, la oración seleccionada es
        introducida por el complementante \forma{que} \exref{ex:06-que}.
  \item Cuando el verbo de la oración principal es interrogativo o dubitativo
        ---por ejemplo, \forma{preguntar, ignorar, querer saber}---, la oración
        seleccionada es introducida por el complementante \forma{si}
        \exref{ex:06-si}.
\end{itemize}
\end{destacado}

El nivel del complementante siempre está presente en la oración, pero, a veces,
puede estar vacío, es decir, puede no haber ninguna marca explícita de la
modalidad en la oración, como normalmente sucede en las oraciones declarativas.

\begin{exe}
  \ex Llegó después de las ocho.
  \ex Ese día Juan estaba feliz.
  \ex\label{ex:06-periodico} Traía un periódico debajo del brazo.
\end{exe}

Otras modalidades oracionales, en cambio, llevan marcadores explícitos. Así, en
\exref{ex:06-talvez}, el adverbio \forma{tal vez} y el modo subjuntivo en el
verbo indican que se trata de una oración dubitativa; en \exref{ex:06-ojala}, el
adverbio \forma{ojalá} y el subjuntivo en el verbo son marcas de la modalidad
desiderativa de la oración; en \exref{ex:06-quetraia}, el índice de la modalidad
interrogativa de la oración es el pronombre interrogativo \forma{qué}, que
representa al constituyente ---\forma{un periódico} de
\exref{ex:06-periodico}--- desplazado en \exref{ex:06-quetraia} a la posición
inicial.

\begin{exe}
  \ex\label{ex:06-talvez} Tal vez llegue después de las ocho.
  \ex\label{ex:06-ojala} Ojalá esté feliz.
  \ex\label{ex:06-quetraia} ¿Qué traía debajo del brazo?
\end{exe}

\begin{recuadro}[Para ampliar: orden en el sintagma de complementante (SC)]
Como se advierte en los ejemplos de arriba, todos los marcadores léxicos
---\forma{tal vez, ojalá, qué}--- se ubican en el margen izquierdo de la
oración, que corresponde al SC. El núcleo del SC contiene la información sobre
la modalidad oracional y allí se coloca la conjunción ---\forma{que, si}---
cuando la oración contiene una subordinada (v. \exref{ex:06-que} y
\exref{ex:06-si}).

En el SC también se alojan otros elementos que sirven para enlazar la oración
con el discurso previo. En este nivel la organización de la oración se vincula
con la \termino{estructura informativa}, que establece la distinción entre tema
y rema, la cual no coincide exactamente con la partición sintáctica en sujeto y
predicado. El \termino{tema} da cuenta de la información ya conocida, previsible
o que el oyente puede inferir; el \termino{rema}, en cambio, aporta la
información nueva.

Si bien no todas las posturas teóricas coinciden, desde la perspectiva que
estamos considerando, un tema destacado en la posición inicial de la oración es
un \termino{tópico}. Como ya adelantamos, los tópicos manifiestan aquello de lo
que se habla en la oración. En \exref{ex:06-top-a}, por ejemplo, el elemento
topicalizado es el COD ---\forma{a Juan}---, que luego es retomado por el
clítico \forma{lo} dentro de la oración. En \exref{ex:06-top-b} y
\exref{ex:06-top-c}, los tópicos aparecen introducidos por marcadores de tópico o
topicalizadores ---\forma{en cuanto a, hablando de}---. Además, como se muestra
en \exref{ex:06-top-c}, el tópico es una categoría recursiva y puede haber más
de uno en la oración.

\begin{exe}
  \ex
    \begin{xlist}
      \ex\label{ex:06-top-a} {}[\etq{SC} [\etq{Top} A Juan\etq{i}] C [\etq{SF/Comp}
          F [\etq{SV/Pred} lo\etq{i} vi esta mañana en la estación]]]
      \ex\label{ex:06-top-b} {}[\etq{SC} [\etq{Top} En cuanto a Juan\etq{i}] C
          [\etq{SF/Comp} F [\etq{SV/Pred} lo\etq{i} vi esta mañana en la
          estación]]]
      \ex\label{ex:06-top-c} {}[\etq{SC} [\etq{Top1} Hablando de Juan\etq{i}],
          [\etq{Top2} a su hermano\etq{j}] C [\etq{SF/Comp} F [\etq{SV/Pred}
          lo\etq{j} vi esta mañana en la estación]]]
    \end{xlist}
\end{exe}

El \termino{foco}, otro de los elementos vinculados con la estructura
informativa incluidos en el SC, representa la información no presupuesta, es
decir, aquella que no forma parte del conocimiento compartido por el hablante y
el oyente. El foco se antepone estableciendo un contraste con las expectativas
ofrecidas por el contexto anterior, por lo que tiene que ser destacado por medio
de recursos enfáticos, como la entonación, \exref{ex:06-foc1} y
\exref{ex:06-foc2}, el orden de palabras \exref{ex:06-foc2}, estructuras de
relieve \exref{ex:06-foc3} o elementos léxicos de focalización, como
\forma{incluso} en \exref{ex:06-foc4}.

\begin{exe}
  \ex\label{ex:06-foc1} {}[\etq{SC} [\etq{Foc} EN LA ESQUINA]\footnote{Convencionalmente,
      la entonación enfática se indica en la escritura por el uso de
      mayúsculas.} fue el accidente] (relieve entonacional)
  \ex\label{ex:06-foc2} {}[\etq{SC} [\etq{Foc} ¿QUÉ] dijiste?] (relieve
      entonacional y orden marcado\footnote{La posición canónica es la esperada
      para un constituyente dentro de la oración, como, por ejemplo, la
      posición inicial para el sujeto o la posverbal para el COD. Cuando los
      elementos se encuentran en la posición canónica, la oración tiene un
      orden no marcado; en cambio, se considera orden marcado cuando se altera
      la disposición de alguno de los constituyentes, como sucede en las
      estructuras tematizadas y focalizadas.} OD -- V)
  \ex\label{ex:06-foc3} En la esquina fue donde ocurrió el accidente
      (estructura de relieve con el verbo \forma{ser})
  \ex\label{ex:06-foc4} {}[\etq{SC} [\etq{Foc} Incluso Juan] estaba enterado del
      asunto] (relieve léxico del foco \forma{Juan} mediante el focalizador
      \forma{incluso})
\end{exe}

También tienen su posición canónica en el margen izquierdo:

\begin{itemize}
  \item Los enlaces con el discurso precedente, como las dos instancias de
        \forma{pero}, que conectan con párrafos anteriores, o el uso inicial de
        \forma{y}, que enlaza con la oración previa, en \exref{ex:06-asimov}:
\end{itemize}

\begin{exe}
  \ex\label{ex:06-asimov}
      % Fragmento literario protegido: se incorpora a mano desde el original.
      \textoaincorporar{el diálogo de I.~Asimov, \textit{Orificios en la
      cabeza} (p.~46), con las dos instancias de \textit{pero} iniciales de
      párrafo y la \textit{y} que abre oración, en cursiva}
\end{exe}

\begin{itemize}
  \item Los \termino{vocativos}, que constituyen una apelación al oyente por
        parte del hablante.
\end{itemize}

\begin{exe}
  \ex \textit{Che}, pero yo realmente podría colaborar en \textit{La Nación} de
      los domingos. (J.~Cortázar, \textit{Rayuela}, p.~222)\pendiente{El
      original dice «Cortázar, R.»}
\end{exe}
\end{recuadro}

\begin{recuadro}[Reflexión final]
Hasta aquí hemos dado un panorama global de la organización oracional
basándonos en los planteos de \textcite{rizzi1995} sobre «la naturaleza en tres
capas o niveles de la oración», a las que hemos vinculado, respectivamente, con
el SV, el SF y el SC. Ahora, en los capítulos que siguen, vamos a adentrarnos en
la predicación, en relación con la cual hay que atender a tres aspectos
interrelacionados, que dan origen a tres estructuras conectadas entre sí:

\begin{itemize}
  \item la \termino{estructura temático-argumental}, que depende de la semántica
        del verbo ---ya sea de acción, de proceso, de estado, etc.--- y, como ya
        hemos anticipado, exige determinada cantidad de argumentos, los cuales,
        según se verá, desempeñan diferentes papeles temáticos (θ) en el evento,
        tales como agente, paciente, experimentante, etc.;
  \item la \termino{estructura sintáctica}, que constituye el enlace entre la
        semántica y la sintaxis, y manifiesta las diferentes funciones
        oracionales que cumplen los papeles temáticos, tales como sujeto, COD,
        COI, etc.;
  \item la \termino{estructura eventiva}, que depende de la naturaleza del
        evento o situación descripta por el verbo ---durativa, puntual, con o
        sin límite final, etc.---, se vincula con la categoría de aspecto, sobre
        todo el léxico (v. capítulos 4 y 8), y está determinada por las distintas
        categorías y constituyentes presentes en la predicación.
\end{itemize}

Asimismo, aunque por el carácter introductorio de este texto no sea posible
explayarnos, cuando la oración se vincula con una situación comunicativa
determinada en la que un emisor se dirige a un receptor, lo que la convierte en
una unidad de realización o enunciado, la manera en que es presentada la
información ---como nueva o conocida, como presupuesta o focalizada---
constituye lo que se conoce como \termino{estructura informativa}, que se
relaciona con los elementos que se ubican en el margen izquierdo de la oración,
correspondiente al nivel del SC.
\end{recuadro}

% --- Secciones finales -------------------------------------------------
\lecturasrecomendadas

\solucionario

\begin{solucion}{\ref{ej:06-determinante}}
Los NN comunes, como \forma{ratas} (1a) y \forma{niños} (1b), por su propia
naturaleza incluyen individuos en una clase ---\forma{Juan es doctor}---. Para
ser usados en referencia a alguna entidad determinada de la realidad
extralingüística deben construirse con un D que los delimite ---\forma{Juan es
el doctor que atiende a mi padre}---, o bien con un Q ---\forma{Llegó algún
amigo de tu padre}---. En cambio, los NN propios, como \forma{Pedro} (2a) o
\forma{Buenos Aires} (2b), tienen la referencia incluida; por eso se pueden usar
sin ningún especificador.
\end{solucion}

\begin{solucion}{\ref{ej:06-oracion}}
b) 2 es una O porque contiene un sujeto y un predicado vinculados por un núcleo
funcional: la F, que manifiesta T (pasado) y Conc (3.ª pers.~sg.). La
particularidad de 2 respecto de 1 es que no contiene un sujeto representado
léxicamente, sino que constituye un caso de lo que se conoce como sujeto tácito.
\end{solucion}

\begin{solucion}{\ref{ej:06-peyrou}}
(3) [\etq{Suj} Su abuelo] contó el caso de uno de ellos\etq{j}, que en una
ocasión [\textit{pro}\etq{[3.ª sg.]}]\etq{j} supo con anticipación la firma de
[\etq{Suj} un decreto] ordenando el pago de cinco años de sueldos atrasados. El
antecedente de \textit{pro} es \forma{uno de ellos}. (4) Inmediatamente
[\textit{pro}\etq{[3.ª sg.]}]\etq{j} visitó durante tres meses, en varias
provincias, las distintas guarniciones. El antecedente de \textit{pro} es
\forma{uno de ellos}, de la oración anterior. (5) A todos los soldados, cabos y
sargentos les [\textit{pro}\etq{[3.ª sg.]}]\etq{j} hizo la misma proposición:
(6) ¿A vos cuánto te [\textit{pro}\etq{[3.ª pl.]}] deben? El antecedente del
primer \textit{pro} es \forma{uno de ellos}; el del segundo \textit{pro} queda
indefinido. (7) Recibida [\etq{Suj} la respuesta],
[\textit{pro}\etq{[3.ª sg.]}]\etq{j} concluía: «[\textit{pro}\etq{[1.ª sg.]}] Te
doy la mitad si me [\textit{pro}\etq{[2.ª sg.]}] firmás recibo por todo». (8)
[\etq{Suj} Los hombres, a veces hambrientos y semidesnudos, y siempre
necesitados] aceptaban, y [\etq{Suj} el hombre] dobló su fortuna en tres meses.
El antecedente del primer \textit{pro} de (7) es \forma{uno de ellos}; el del
segundo \textit{pro} es el emisor; el del tercer \textit{pro}, uno de los
receptores.\pendiente{La numeración de las oraciones del solucionario no
coincide con la del texto; el original usa el índice \textit{i} para
\textit{uno de ellos}}
\end{solucion}

\begin{solucion}{\ref{ej:06-esoracion}}
3. O formada por el SN sujeto \forma{un meteorito} y el SV predicado
\forma{cayó en la autopista}. 4. O formada por un sujeto tácito y el SV
predicado \forma{vendrán más tarde}. 5. SAdv. 6. O formada por el SN sujeto
\forma{solo ranchos y casuchas con techo de paja} y el SAdv predicado
\forma{allá}. 7. O formada por un SV de V impersonal: \forma{En el cuento había
un ogro, un burro y una princesa}. 8. Hay dos oraciones coordinadas: la primera
está formada por el SN sujeto \forma{los familiares cercanos} y el SV predicado
\forma{primero entraron}; la segunda está formada por el SN sujeto \forma{los
demás parientes} y un SV con el V elidido, que puede reponerse porque es el mismo
que el de la primera: \forma{después (entraron)}. 9. O formada por un SV de V
impersonal de tipo meteorológico: \forma{Granizó copiosamente en varias zonas
del interior de la provincia}. 10. O formada por el SN sujeto \forma{el café de
la sobremesa} y el SA predicado \forma{infaltable}.

\textit{Aclaración}: los ejercicios que siguen son preparatorios para los del
capítulo siguiente, que ya especifican las diferencias entre argumentos. Aquí
solo se solicita identificar los argumentos y las funciones sintácticas que
cumplen.
\end{solucion}

\begin{solucion}{\ref{ej:06-argumentos}}
Respuestas posibles. 2. a) [\etq{Suj} La madre de Lucía] prepara [\etq{COD}
unas exquisitas tortas de chocolate]. b) [\etq{Suj} Algunos aspirantes]
aprobaron [\etq{COD} el examen]. 3. a) [\etq{Suj} El gerente de marketing]
presentó [\etq{COD} un informe] [\etq{COI} al director de la compañía]. b)
[\etq{Suj} Los organizadores del evento] enviaron [\etq{COD} invitaciones]
[\etq{COI} a los conferencistas].
\end{solucion}

\begin{solucion}{\ref{ej:06-funciones}}
a) 2. \forma{Informó} exige tres argumentos, que en este caso son \forma{el
primer ministro, la población} y \forma{las nuevas medidas de gobierno}. 3.
\forma{Guardar} exige tres argumentos: \forma{el fiscal del caso, los
expedientes} y \forma{en un sobre}. b) 2. [\etq{Suj} El primer ministro] informó
[\etq{COI} a la población] [\etq{COD} las nuevas medidas de gobierno]. 3.
[\etq{Suj} El fiscal del caso] guardó [\etq{COD} los expedientes] [\etq{CCL} en
un sobre]. Aclaración: \forma{en un sobre} es un argumento exigido por la
semántica del V, que sin su presencia quedaría incompleto
(\juicio{*}{\forma{El fiscal del caso puso los informes}}); por lo tanto,
se considera un argumento del verbo con la función sintáctica de complemento
circunstancial de lugar.\pendiente{En la prueba, el original pasa a «puso los informes»; en 2 el análisis omite \textit{sobre} ante \textit{las nuevas
medidas}}
\end{solucion}

\begin{solucion}{\ref{ej:06-llover}}
b) Respuestas posibles: \forma{Granizó muy fuertemente esta mañana. Tronaba
tanto que me asusté}. c) Dado que en inglés o francés no es posible omitir el
sujeto, los VV meteorológicos se construyen con un sujeto que se pronuncia, pero
no tiene significado. Ejemplos: inglés \forma{It snows} (\sig{nieva}); francés
\forma{Il neige}. d) Este elemento que «hace las veces de sujeto» se conoce como
\termino{expletivo}. e) La exigencia de un expletivo se relaciona con el
parámetro del sujeto, que en estas lenguas obligatoriamente debe estar expreso
en la oración, a diferencia del español o el italiano, que permiten la omisión.
Ejemplos: español \forma{Nieva}; italiano \forma{Nevica}. f) Las lenguas como el
español o el italiano, que permiten la omisión del sujeto, se conocen como
\textit{pro-drop}.
\end{solucion}

\begin{solucion}{\ref{ej:06-cuatro}}
a) 2. [María] le compró [su auto] [a su prima] [en 500\,000\,\$].
\textit{Aclaración}: \forma{le} y \forma{a su prima} son correferenciales y
expresan el mismo argumento. b) Respuestas posibles: [El profesor de Química]
viene todos los días [a la Facultad] [desde Tigre] [por la Panamericana]. / [Los
exportadores de soja] vendieron [el cereal] [al exterior] [a precio récord] este
año.
\end{solucion}
